\setcounter{section}{1}
\setcounter{theorem}{3}
\section{Notations and Preliminaries}
All logarithms in the paper are base 2.

We use some standard graph theory terminology: If $G$ is a directed graph and $(u,v)$ is an edge, $v$ is called the \emph{head} of the edge and $u$ the \emph{tail}. Our directed graphs are always acyclic with designated source vertex $s$ and sink vertex $t$. The \emph{depth} of a vertex $v$, denoted $\depth(v)$, is the length (in edges) of a longest path from $s$ to $v$. The depth of the graph, denoted by $\depth(G)$, is the depth of $t$.

For any two vertices $u$ and $v$ in an ABP, the polynomial computed between $u$ and $v$ is the sum of weights of all paths between $u$ and $v$  in the ABP. We denote this by $[u,v]$. 

The formal degree of a vertex $u$ in an ABP denoted $\fdeg(u)$, is defined inductively as follows: If $s$ is the start vertex of the ABP, $\fdeg(s)=0$. If $u$ is a vertex with incoming edges from $u_1, \ldots, u_k$, labeled by non-zero polynomials $\ell_1, \ldots, \ell_k$, respectively, then
\[
\fdeg(u) = \max_{i \in [k]} \set{ \deg(\ell_i) + \fdeg(u_i) }.
\]
It follows by induction that for every vertex $u$, $\deg ([s,u]) \le \fdeg(u)$ (however, cancellations can allow for arbitrary gaps between the two). Also, note that the formal degree of vertices is monotonically non-decreasing along any path from the source vertex to the sink vertex. The formal degree of the ABP is the maximal formal degree of any vertex in it.

We sometimes denote by $\vx$ the vector of variables $(x_1, \ldots, x_n)$, where $n$ is understood from the context. Similarly we use $\mathbf{0}$ to denote the $n$-dimensional vector $(0,0,\ldots,0)$.

	

\subsection{A decomposition lemma}
The following lemma gives a decomposition of a (possibly unlayered) ABP in terms of the intermediate polynomials it computes. Its proof closely resembles that of Lemma 3.5 of \cite{K19}. For completeness we prove it here for a slightly more general model.

\begin{lemma}[\cite{K19}]\label{lem:structure}
	Let $\mathcal{B}$ be a (possibly unlayered) algebraic branching program which computes a degree $d$ polynomial $P\in\F[x_1,\ldots,x_n]$, and has formal degree $d$. Further, assume that the edges of $\mathcal{B}$ are labelled by arbitrary polynomials of degree at most $\Delta$, where $1\le\Delta\le d/2$. Set $d'=\floor{d/\Delta}$.
		
	 For any $i \in \set{1, 2, \ldots ,d'-1}$, let $S_i = \set{u_{i,1}, u_{i,2}, \ldots, u_{i,m}}$ be the set of all vertices in $\mathcal{B}$ whose formal degree is in the interval $[i \Delta, (i+1)\Delta)$.
	 
	 Then, there exist polynomials $Q_{i,1}, Q_{i,2}, \ldots, Q_{i,m}$ and $R_i$, each of degree at most $d-1$ such that
		\[P = \sum_{j=1}^{m} [s, u_{i,j}] \cdot Q_{i,j} + R_i\, .\]
\end{lemma}

\begin{proof}
Fix $i$ as above and set $S_i = \set{u_{i,1}, u_{i,2}, \ldots, u_{i,m}}$ as above (observe that since each edge label is of degree at most $\Delta$, $S_i$ is non empty). Further suppose, without loss of generality, that the elements of $S_i$ are ordered such that there is no directed path from $u_{i,j}$ to $u_{i,j'}$ for $j' > j$.

Consider the unlayered ABP $\mathcal{B}_1$ obtained from $\mathcal{B}$ by erasing all incoming edges to $u_{i,1}$, and multiplying all the labels of the outgoing edges from $u_{i,1}$ by a new variable $y_1$. The ABP $\mathcal{B}_1$ now computes a polynomial of the form
\[
	P'(y_1, x_1, \ldots, x_n) = y_1 \cdot Q_{i,1} + R_{i,1}
\]
where $P = P'([s,u_{i,1}], x_1, \ldots, x_n)$. $R_{i,1}$ is the polynomial obtained from $\mathcal{B}_1$ by setting $y_1$ to zero, or equivalently, removing $u_{i,1}$ and all its outgoing edges. We continue in the same manner with $u_{i,2}, \ldots, u_{i,m}$ to obtain
\[
P = \sum_{j=1}^m [s,u_{i,j}] \cdot Q_{i,j} + R_i.
\]
Indeed, observe that since there is no path from $u_{i,j}$ to $u_{i,j'}$ for $j' > j$, removing $u_{i,j}$ does not change $[s,u_{i,j'}]$. The bound on the degrees of $Q_{i,j}$ is immediate from the fact that the formal degree of the ABP is at most $d$ and $\fdeg(u_{i,j}) \ge 1$. It remains to argue the $\deg(R_i) \le d-1$.

The polynomial $R_i$ is obtained from $\mathcal{B}$ by erasing all the vertices in $S_i$ and the edges touching them. We will show that every path in the corresponding ABP computes a polynomial of degree at most $d-1$. Let $s=v_1, v_2, \ldots, v_r=t$  be such a path, which is also a path in $\mathcal{B}$. Let $v_k$ be the minimal vertex in the path whose degree (in $\mathcal{B}$) is at least $(i+1)\Delta$
(if no such $v_k$ exists, the proposition follows).
As $v_{k-1} \not \in S_i$, the formal degree of $v_{k-1}$ is at most $i \Delta - 1$.
The degree of the polynomial computed by this path is thus at most $i\Delta - 1 + \Delta + D = (i+1)\Delta - 1 + D$, where $D$ is the degree of product of the labels on the path $v_k, v_{k+1}, \ldots, t$. To complete the proof, it remains to be shown that $D \le d - (i+1)\Delta$.

Indeed, if $D \ge d - (i+1)\Delta + 1$ then since the degree of $v_k$ is at least $(i+1)\Delta$, there would be in $\mathcal{B}$ a path of formal degree at least $(i+1)\Delta +D \ge d + 1$, contradicting the assumption on $\mathcal{B}$.
\end{proof}


	
\subsection{Variety and its Dimension}
One of the important notions we will use in our proofs is that of an affine algebraic variety. Given a set of polynomials, the algebraic variety defined by these polynomials is defined to be the set of their common zeros. That is, if $S$ is a set of polynomials on $n$ variables over a field $\F$, then
\[\mathcal V=\{\mathbf a\in\F^n:\forall f\in S,f(\mathbf a)=0\}.\]
Given a variety, an important property that is studied is its dimension. Intuitively, it is an appropriate generalisation of the notion of dimension for linear spaces. We will not be defining it formally here and refer the interested reader to the book by Cox, Little and O'Shea \cite{CLO}. However we state formally the properties, of dimensions of varieties, that we will be using in our proofs.
\begin{lemma}[Section 2.8 in {\cite{smith2014introduction}}]\label{lem:dimlower}
Let $S$ be a set of polynomials in $n$ variables over an algebraically closed field $\F$ such that $|S|\le n$. Let $\mathcal V=\mathbb V(S)$ be the set of common zeros of polynomials in $S$. If $\mathcal V$ is non-empty, then the dimension of $\mathcal V$ is at least $n-|S|$.
\end{lemma}
\begin{lemma}[Chapter 4 in {\cite{CLO}}]\label{lem:diminclusion}
Let $\F$ be an algebraically closed field, and let $\mathcal V_1\subseteq\F^n$ and $\mathcal V_2\subseteq\F^n$ be two affine varieties such that $\mathcal V_1\subseteq\mathcal V_2$. Then, the dimension of $\mathcal V_1$ is at most the dimension of $\mathcal V_2$.
\end{lemma}
